\section{图搜索、最短路和状态建模}

图题不是看到网格就写 DFS，也不是看到最短就写 BFS。第一步永远是建模：节点是什么，边是什么，边权是什么，状态是否还包含钥匙、剩余障碍次数、访问集合或时间。岛屿数量的节点是陆地格子，边是四方向相邻；课程表的节点是课程，边是先修依赖；开锁问题的节点是四位密码，边是转动一位；最小体力路径的节点是格子，边权是高度差，但路径代价是最大边权。

BFS 的正确性来自层次扩展。无权图中，从起点出发，队列先处理距离为零的状态，再处理距离为一的状态，然后是距离为二。只要访问标记在入队时设置，一个状态第一次入队时的距离就是最短距离。若等到出队才标记，同一状态可能被同层多个前驱重复加入，虽然有时答案仍对，复杂度却会膨胀。多源 BFS 只是把所有源点同时放进距离为零的层。

DFS 更适合连通块、路径搜索和状态检测。连通块题中，一次 DFS 或显式栈搜索会标记从起点可达的全部节点；拓扑判环中，三色标记能区分未访问、当前路径中、已完成；回溯搜索中，DFS 路径状态可以在进入和退出时维护。递归写法短，但工程上要认识深度风险，尤其是网格很大或树退化成链时。

拓扑排序是有向图依赖关系的核心工具。入度为零的节点表示当前没有未满足依赖，可以安全输出；输出它以后，只会减少它指向的后继入度。若最终输出节点数少于总数，剩余节点一定在环中或被环依赖。边方向必须讲清楚：课程表中，先修课指向后续课程，还是课程指向先修课，决定输出顺序是否正确。

Dijkstra 适用于非负边权。它维护当前已知最短距离，每次从堆中取出距离最小的未过期状态。非负边保证从这个状态再走出去不会产生一条更短路径回头修改它。C++ 的 priority\_queue 不支持 decrease-key，所以常用做法是更新距离时压入新状态，弹出时若距离不是当前 distance 数组中的值，就说明它过期，直接跳过。

有些题看起来不像普通最短路，但仍满足 Dijkstra 的单调扩展。最小体力路径中，路径代价不是边权和，而是路径上最大高度差；扩展一条边后，新代价是当前代价和边代价的最大值。这个值不会随着路径延长而降低，因此仍可用堆按当前最小代价扩展。另一种解法是二分体力阈值，再用 BFS 检查是否可达，因为阈值越大，可走边越多。

状态压缩 BFS 要特别小心 visited 的维度。位置相同但钥匙集合不同，未来能力不同；位置相同但剩余障碍消除次数不同，未来能力也不同。此时 visited 不能只按位置记录，而要按位置加资源状态记录。很多看似超时或错误的 BFS，根本原因是把不同状态错误合并，导致要么漏解，要么重复搜索。

实验建议：把腐烂橘子写成每分钟全图扫描和多源 BFS 两版，比较重复扫描次数。把网络延迟时间写成普通 BFS 和 Dijkstra 两版，用不同权重反例证明 BFS 错误。再把带钥匙迷宫的 visited 从二维改成三维，观察二维 visited 如何错误剪掉合法路径。

\begin{principlebox}{本节复盘}
阅读本节后，请回到相邻章节的例题，把暴力瓶颈、优化依据、状态不变量、C++ 容器成本和边界测试重新写一遍。能把同一思想迁移到变体题，才算真正掌握。
\end{principlebox}
